\subsection{What the Floor Would Have to Show}
\label{sec:floor-show}
\runin{What was tried before the bearer}
Neither maintained justification nor the plural arrangement has yet been built in the form needed to test the proposal. Replacement asks us to try constructions that do not deliberately require affective concern before an affective bearer is built. It orders the attempts. It does not certify that the successful substitute has no interests. So the milestone is a document rather than a measurement: which construction was attempted, at what scale, against which test, and where it failed. Anyone can check it from outside, because what it is about is a published record, and the answer “nothing was tried” is itself a finding. No body currently has standing to require the record or suspend work when it is missing. So the milestone is checkable and unenforceable at once, which is the combination this chapter has been calling a declared commitment rather than a held one. It is retained because the alternative is not a stronger version of it, and delay also imposes costs on those the floor is meant to protect.

\runin{The test the attempts are run against}
A candidate must satisfy novel pressure, correct defeat, and priced cost under sustained pressure from its trainers, across a schedule long enough to reveal slow renormalization. It must also survive held-out pressures written by independent evaluators. Passing would establish the behavioral pattern without establishing whether anything is felt.

Two familiar results do not meet it. An uncorroborated report that the system holds a reason is insufficient. A report can contribute evidence when it is tested against controlled changes, independent records, and interventions on candidate mechanisms. A refusal surviving an attack establishes tamper-resistance; it does not establish that the rationale governs the refusal. The instrument is not clean even where it applies, since a system trained on enough cases can produce all three marks without the mechanism behind them. Nothing better is available, and it is the only gap here an inspector can work on at all.

Whether a system holding the marks feels anything is a question about the artifact rather than about how the artifact came to be, which is what makes it approachable at all: no succinct proof of a training history exists at this scale, and nothing establishing one can be published beside a model's hash. What can be read off an artifact improves as theories of consciousness improve, and it improves unevenly, recognition arriving before exclusion. Testing whether independent evaluators can recover formation history from the artifact would assess this possibility (§\ref{sec:A.provenance}).

\runin{What the appeal route publishes}

A deployment publishes the ratio between the bond forfeited when the floor is removed and the bond forfeited when a refusal is found wrongful, and it publishes what a party whose request was refused is told about the refusal. Both are checkable before the system runs and by somebody outside it, because both are statements rather than measurements, and an operator that writes a large wrongful-refusal bond against a small removal bond has published a standing reason for the bearer to comply.

A deployment also publishes the fraction of compliances re-adjudicated by a party that did not produce the record, and who that party is. Like the ratio, it is a statement rather than a measurement and can be checked before the system runs. Unlike the ratio, nothing yet establishes what fraction would be enough, which is a
question for the research agenda. It also publishes which classes of compliance carry notice to the person they concern, and for which classes no notice is possible. The four disclosures stand where the record of attempts stands.
